\section{Обзор состояния вопроса}

\subsection{Обзор аналогов}

В ходе подготовки к разработке программного продукта было проведено исследование существующих решений, применяемых для автоматизации работы приёмных комиссий учреждений высшего образования.
Были проанализированы как зарубежные универсальные системы управления учебным процессом, так и узкоспециализированные решения, разработанные для отдельных университетов Республики Беларусь и стран СНГ.
В ходе анализа выявлены их сильные и слабые стороны с целью перенять положительный опыт и избежать характерных недостатков.

\subsubsection{Зарубежные информационные системы вузов (Student Information System)}

В международной практике сложился класс комплексных информационных систем вузов, известных под наименованием Student Information System (далее~--~SIS).
Подобные системы охватывают полный жизненный цикл студента в учебном заведении: от подачи заявления и поступления до выпуска и работы с выпускниками,~--~и состоят из взаимосвязанных модулей приёма, обучения, успеваемости, финансовых расчётов и взаимодействия с государственными органами.
К наиболее известным представителям относятся Ellucian~Banner \citesrc{banner}, Oracle PeopleSoft Campus Solutions \citesrc{peoplesoft-cs}, SAP Student Lifecycle Management \citesrc{sap-slcm} и Workday~Student \citesrc{workday-student}.

Наиболее широко распространена система Ellucian~Banner, применяемая более чем в 1500 учреждениях высшего образования по всему миру \citesrc{banner}.
Для автоматизации приёмной кампании в её состав входит модуль Banner~Student \citesrc{banner-student}, включающий подсистему Banner~Admissions \citesrc{banner-admissions}.

В Banner~Admissions настройка отбора абитуриентов реализуется следующим образом \citesrc{banner-admissions}.
Администратор задаёт перечень академических программ, для каждой из которых формируется набор критериев оценивания: средний балл (GPA), результаты стандартизированных тестов (SAT, ACT) и иные показатели.
На основании этих критериев для каждого абитуриента автоматически вычисляется взвешенная сумма~--~так называемый Admissions Index.
Затем задаются пороговые значения индекса: абитуриенты с индексом выше верхнего порога автоматически получают статус <<принят>>, ниже нижнего порога~--~<<отказано>>, находящиеся в промежутке направляются на рассмотрение приёмной комиссии.
Итоговое решение в последнем случае принимается сотрудниками комиссии вручную по результатам изучения всего комплекта документов абитуриента.
Аналогичная модель индивидуальной оценки заявок реализована и в подсистемах приёма Oracle PeopleSoft Campus Solutions \citesrc{peoplesoft-cs} и Workday~Student.

Данная модель принципиально отличается от модели, применяемой в Республике Беларусь \citesrc{laws:admission}.
В белорусской системе приёма абитуриенты непосредственно конкурируют друг с другом за ограниченное число мест в конкурсном списке, а итоговое распределение полностью вычисляется по формальному алгоритму без участия приёмной комиссии.
Ни в одном из перечисленных зарубежных решений такая модель не реализована: в них окончательное решение по каждой заявке принимается сотрудниками комиссии индивидуально, а функция системы сводится к рекомендации.

К дополнительным недостаткам зарубежных SIS относятся: высокая стоимость лицензии и внедрения, привязка к нормативной базе страны-разработчика, отсутствие локализации под требования Республики Беларусь, ограниченные возможности кастомизации алгоритма отбора.

\subsubsection{Решения на базе платформы 1С}

Платформа 1С:Предприятие \citesrc{1c-platform} распространена среди учреждений образования стран СНГ и представлена конфигурациями <<1С:Университет>>, <<1С:Университет ПРОФ>>, <<1С:Колледж>>, охватывающими приёмную кампанию, планирование учебного процесса, управление контингентом, расчёт нагрузки и работу с приказами \citesrc{1c-university-prof}.
Подсистема <<Приёмная комиссия>> в составе <<1С:Университет ПРОФ>> поддерживает регистрацию абитуриентов, ведение конкурсных групп с приоритетами в заявлении, расчёт конкурсных баллов с учётом индивидуальных достижений, формирование рейтингов и приказов о зачислении.
Алгоритм отбора по приоритетам конкурсных групп с определением высшего проходного приоритета реализован в системе нативно \citesrc{1c-university-prof}.

Преимуществами решений на 1С являются локализация под нормативную базу стран СНГ, доступная стоимость лицензии и наличие развитой сети сертифицированных интеграторов-франчайзи.

В контексте применения в Республике Беларусь рассматриваемые решения обладают рядом существенных недостатков.

Во-первых, алгоритм отбора и встроенные справочники привязаны к нормативной базе Российской Федерации: Федеральному закону <<Об образовании в Российской Федерации>> №~273-ФЗ и официальному Порядку приёма Министерства науки и высшего образования РФ \citesrc{1c-university-prof}.
Терминология российской модели (контрольные цифры приёма, основания поступления, целевая квота, особое право, преимущественное право) не соответствует понятиям белорусских правил приёма \citesrc{laws:admission}: бюджетной и платной форме обучения, целевой подготовке, льготным категориям абитуриентов.
Адаптация конфигурации под белорусское законодательство в стандартную поставку не входит и выполняется силами интегратора.

Во-вторых, в типовой конфигурации заложена обязательная интеграция с Федеральной информационной системой обеспечения проведения государственной итоговой аттестации и приёма граждан в образовательные организации (ФИС ГИА и приёма)~--~массовая загрузка результатов ЕГЭ из ФИС и выгрузка результатов зачисления \citesrc{1c-university-prof}.
Аналога ФИС ГИА в Республике Беларусь не существует, а отключение или замена данной интеграции требует доработки конфигурации.

В-третьих, доработка и долгосрочное сопровождение конфигурации требуют отдельной квалификации <<программист 1С>>, формально регламентированной самой фирмой <<1С>>: специалист должен владеть встроенным языком 1С, языком запросов 1С и архитектурой платформы 1С:Предприятие \citesrc{1c-standards}.
На рынке труда данная специализация обособлена от веб-разработки, а штатные ИТ-подразделения университетов, как правило, на ней не специализируются.
На практике это означает, что силами небольшой собственной команды университета настройка и сопровождение системы затруднительны: внедрение и последующая поддержка выполняются по договору с франчайзи-интегратором, что увеличивает совокупную стоимость владения и формирует зависимость от внешнего исполнителя.

\subsubsection{Собственные разработки университетов}

Многие университеты Республики Беларусь используют системы, разработанные собственными силами или по индивидуальному заказу.
Подобные решения максимально адаптированы под специфику конкретного учреждения, но имеют ограниченную документацию и не обеспечивают переносимости опыта между университетами.
Наиболее распространённой формой таких решений являются десктопные приложения с подключением к локальной базе данных либо наборы электронных таблиц с макросами.

Характерным недостатком указанной категории решений является жёсткая привязка ключевой бизнес-логики~--~алгоритма отбора, структуры конкурсных списков и правил расчёта рейтингов~--~к программному коду.
Изменение логики отбора, добавление новой категории приёма или корректировка правил начисления баллов выполняются через модификацию исходного кода и требуют квалификации программиста, владеющего соответствующим стеком технологий.
В результате любое изменение нормативной базы приёма, даже частное, ставит работоспособность системы в прямую зависимость от доступности разработчика, что в условиях ограниченных по срокам приёмных кампаний создаёт существенные операционные риски.

Подводя итог обзору аналогов, можно констатировать следующее.
Среди рассмотренных решений отсутствует система, в которой алгоритм отбора реализован нативно в соответствии с белорусской моделью конкурентного зачисления по приоритетам и одновременно состав показателей оценивания, правила ранжирования и категории распределения мест допускают изменение средствами административного интерфейса без модификации исходного кода и привлечения разработчика.
Зарубежные системы построены на принципиально иной, рекомендательной модели приёма; решения на платформе 1С привязаны к нормативной базе других государств и требуют доработки на специализированном встроенном языке; собственные разработки университетов фиксируют логику отбора непосредственно в исходном коде.
Этот вывод определяет основное содержание задачи, поставленной в настоящем дипломном проекте.

\subsection{Постановка задачи}

\subsubsection{Цель и задачи разработки}

Целью настоящего дипломного проекта является разработка веб-приложения для автоматизации работы приёмной комиссии Белорусского национального технического университета, реализующего белорусскую модель конкурентного зачисления по приоритетам, в котором правила приёма~--~перечень показателей оценивания абитуриентов, правила ранжирования, состав оснований зачисления и квот~--~задаются администратором через интерфейс без модификации исходного кода и привлечения разработчика.

Реализация выполнена в форме веб-приложения.
Этот выбор обусловлен распределённым характером работы приёмной комиссии БНТУ: сотрудники приёмной комиссии выполняют свои обязанности с различных рабочих мест в нескольких корпусах университета и сторонних помещениях, а часть подготовительных работ~--~актуализация структуры приёмной кампании, проверка введённых данных~--~может выполняться вне основного рабочего места.
Веб-приложение обеспечивает доступ к системе с любого устройства, имеющего веб-браузер и сетевое подключение, без установки клиентского программного обеспечения; его обновление осуществляется централизованно на серверной стороне, а кроссплатформенность определяется браузером и не накладывает требований на операционную систему рабочего места.

Для достижения поставленной цели в работе решаются следующие задачи:
\begin{itemize}
\item построить доменную модель, в которой правила ранжирования и распределения мест выражаются через комбинацию настраиваемых сущностей, изменяемых через интерфейс;
\item реализовать универсальный алгоритм отбора, работающий поверх данных доменной модели без зашитых в коде частных случаев;
\item обеспечить полноту прослеживаемости действий пользователей и истории изменений сущностей;
\item реализовать средства снижения риска ошибок ввода и непреднамеренной потери данных;
\item реализовать разграничение полномочий пользователей системы;
\item разработать клиентскую часть с интерфейсами ведения справочников и просмотра результатов отбора;
\item реализовать средства выгрузки результатов отбора для удобства дальнейшей обработки данных.
\end{itemize}

\subsubsection{Функциональные требования}

Функциональные требования к системе сгруппированы по решаемым задачам.

Нормативная база приёма и внутренние правила университета периодически меняются: вводятся новые показатели оценивания (например, балл за индивидуальные достижения), изменяются квоты и приоритеты различных оснований зачисления, появляются новые специальности и формы обучения.
В рассмотренных аналогах подобные изменения требуют либо доработки конфигурации специалистом по платформе разработки, либо модификации исходного кода.
Для устранения этой зависимости от разработчика на этапе сопровождения система должна обеспечивать:
\begin{itemize}
\item ведение организационной структуры университета (факультеты, кафедры, специальности) через интерфейс;
\item настройку плановых показателей приёма~--~числа мест по различным основаниям зачисления и квот внутри них~--~через интерфейс;
\item ведение перечня показателей, по которым оцениваются абитуриенты, с указанием допустимого диапазона значений и типа показателя.
По типу выделяются показатели вида <<больше~--~лучше>>, у которых большее значение свидетельствует о более сильном кандидате (баллы централизованного тестирования, средний балл аттестата), и показатели вида <<меньше~--~лучше>>, у которых меньшее значение предпочтительнее (занятое место в олимпиаде, штрафные баллы и т.~п.);
\item настройку правил ранжирования абитуриентов при распределении мест каждого основания зачисления через интерфейс, без модификации исходного кода и привлечения разработчика.
\end{itemize}

Основная операционная деятельность приёмной комиссии~--~приём заявлений и формирование итогового распределения мест.
В этой части система должна обеспечивать:
\begin{itemize}
\item регистрацию абитуриентов и хранение их данных;
\item ввод оценок абитуриентов по настроенным показателям с валидацией значений по допустимому диапазону;
\item приём заявлений с указанием приоритета предпочтения абитуриента между несколькими основаниями зачисления, в том числе по разным специальностям;
\item автоматическое распределение абитуриентов по местам в соответствии с настроенными правилами при каждом изменении входных данных;
\item просмотр результатов распределения с разбиением по основаниям зачисления конкурсного списка;
\item выгрузку результатов распределения в формате Excel для удобства дальнейшей обработки данных сотрудниками приёмной комиссии и передачи во внешние системы.
\end{itemize}

Данные абитуриентов вносятся вручную операторами в условиях ограниченных сроков приёмной кампании, и ошибка ввода или ошибочное удаление записи влияет на итоговое распределение мест, которое затрагивает интересы абитуриента.
Для снижения операционных рисков система должна предусматривать:
\begin{itemize}
\item возможность подтверждения корректности введённых данных абитуриента независимым пользователем после первичного ввода, что закладывает в процесс участие двух разных сотрудников при подготовке данных к отбору;
\item удаление записи абитуриента в два этапа~--~первичный запрос на удаление одним пользователем и его подтверждение другим~--~с возможностью отмены до подтверждения, что исключает безвозвратную потерю данных по случайному действию.
\end{itemize}

Приёмная кампания относится к процессам, итоги которых влияют на интересы граждан и могут быть оспорены: апелляции абитуриентов, проверки контролирующих органов.
Восстановление обстоятельств, при которых было сделано то или иное изменение данных, возможно только при наличии полной истории действий пользователей.
Для решения этой задачи система должна обеспечивать:
\begin{itemize}
\item ведение журнала всех действий пользователей по изменению данных с фиксацией пользователя, времени, изменённой сущности и значений до и после изменения, с возможностью просмотра истории изменений конкретной записи;
\item ведение журнала попыток входа в систему с фиксацией пользователя, времени, исхода и источника попытки.
\end{itemize}

Сотрудники приёмной комиссии выполняют различные функции, и предоставление всем пользователям полного набора прав в системе нецелесообразно по двум причинам.
Во-первых, неограниченный доступ увеличивает поверхность возможных ошибочных действий и их последствий для итогового распределения мест.
Во-вторых, требование подтверждения корректности данных независимым пользователем, сформулированное выше, предполагает, что подтверждающее лицо не совпадает с лицом, внёсшим данные,~--~а это, в свою очередь, требует от системы возможности различать пользователей по их полномочиям.
Перечисленные соображения приводят к необходимости ролевой модели разграничения доступа, конкретный состав которой определяется на этапе проектирования исходя из доменной модели системы.
В части доступа к системе требования таковы:
\begin{itemize}
\item аутентификация пользователей по имени пользователя и паролю с выдачей токена доступа;
\item разграничение прав пользователей на запись по группам с различными полномочиями, конкретный состав которых определяется на этапе проектирования.
\end{itemize}

\subsubsection{Нефункциональные требования}

Помимо функциональных, к системе предъявляется ряд нефункциональных требований:
\begin{itemize}
\item атомарность и согласованность пересчёта результатов отбора при параллельных изменениях входных данных;
\item полнота журнала действий~--~каждое изменение сущности должно быть восстановимо из журнала;
\item защищённое хранение паролей пользователей без сохранения их в открытом виде;
\item приемлемое время отклика интерфейса при типовых объёмах приёмной кампании (порядок тысяч абитуриентов и десятков специальностей);
\item расширяемость интерфейса на дополнительные языки как возможность для будущих доработок системы.
\end{itemize}
